\begin{graphicspathcontext}{{./chapters/ai/imgs/},{./chapters/ai/imgs/auto/},\old}

\sidecite{Brynjolfsson2014,Acemoglu2018,Freedman2018,Zuboff2019}
\begin{frame}[t]{{Impacts économiques :} emploi, productivité, inégalités}
	\vspace{-.25cm}
	\begin{columns}[T]
		\begin{column}{.5\linewidth}
			\begin{block}{Opportunités}
				\begin{itemize}
				\item Prise en charge de tâches à risques
				\item Aide à la décision
				\item Gains de productivité%\cite{Brynjolfsson2014}
				\item Nouvelles tâches à valeur ajoutée
				\end{itemize}
			\end{block}
		\end{column}
		\begin{column}{.5\linewidth}
			\begin{block}{Risques}
				\begin{description}
				\item[Suppression d'emplois] tâches répétitives, certains métiers intellectuels%\cite{Acemoglu2018}
				\item[Concentration des richesses] grandes entreprises maîtrisant l'IA%\cite{Zuboff2019}
				\item[Précarisation des travailleurs] surveillance algorithmique, micro-travail%\cite{Freedman2018}
				\end{description}
			\end{block}
		\end{column}
	\end{columns}
	\vspace{.5cm}
	\alertbox{\Emph{Question centrale} : Comment partager les bénéfices de l'IA sans creuser les inégalités ?}
	\alertbox*{\emph{Pistes :} formation continue, fiscalité spécifique, revenu universel...}
\end{frame}

\end{graphicspathcontext}

\endinput
